% New finite-proper descent consequence; awaiting independent actual-file audits.
\begin{corollary}[Proved: bounded-centre projective growth after finite normalization]
\label{D:nonnormal-bounded-centre-projective-growth}
Let $S$ be a proper scheme of dimension at most two over an algebraically
closed field of characteristic $p>0$. Let $T$ be a finite closed subset
of codimension at least two in every irreducible component it meets,
and let $U$ be a connected component of $S\setminus T$. Fix an integer
$B\ge1$. There is a constant $C_{S,T,B}$ such that, for $q=p^f$ and
$r\ge2$ with $\gcd(r,q-1)\le B$,
\[
 \#\bigl(\operatorname{Hom}_{\mathrm{cont}}
 (\pi_1(U),\operatorname{PSL}_r(\mathbf F_q)) /
 \operatorname{PGL}_r(\mathbf F_q)\bigr)
 \le q^{C_{S,T,B}r^2}.
\]
Consequently the quotient formed from all distinct nonabelian simple
quotients of these types is topologically finitely generated. The
ranks are unrestricted, while the natural central order is bounded
by the fixed $B$.
\end{corollary}
\begin{proof}
Reduction does not change finite etale covers. Normalize the reduced
proper scheme. Its normalization is finite, and its inverse image
$U'\to U$ is a finite surjection with finitely many connected
components $U_i$. Each two-dimensional component is a normal proper
surface minus finitely many points. Proposition~\ref{D:bounded-centre-projective-growth}, applied by
the same finite-cover argument as
Corollary~\ref{D:surface-bounded-centre-simple-product}, gives a bound
$q^{c_{i,B}r^2}$ for its projective conjugacy classes. The components
of dimension at most one are proper: the codimension hypothesis
prevents $T$ from meeting those components. Their fundamental groups
are finitely generated, and hence their representations into any
finite group $P$ number at most $|P|^{d_i}$ for a fixed $d_i$.
We apply this with $P=\operatorname{PGL}_r(\mathbf F_q)$ and use
$|P|\le q^{r^2}$.

We give the nonabelian descent count explicitly. A homomorphism to
$\operatorname{PSL}_r(\mathbf F_q)$ determines a torsor for the
constant finite group $P$ by extension of the structure group.
Two such homomorphisms determine isomorphic $P$-torsors precisely
when they are conjugate by $P$. Thus the classes which we wish to
count inject into the set of $P$-torsor classes on $U$. Their
restrictions to $U'$ range over at most
$q^{(\sum_i c_{i,B})r^2}$ tuples, after absorbing the curve and
point bounds. This use of $P$-torsors avoids replacing projective
conjugacy by conjugacy inside $\operatorname{PSL}_r$.

Fix one such tuple of $P$-torsors. The two pullbacks to
$U'\times_U U'$ have a finite etale Isom sheaf whose geometric
fibres, when nonempty, have cardinality $|P|$. A section on a
connected component is determined by its value at one geometric
point: the equalizer of two sections of a finite etale separated
sheaf is open and closed. If $c$ is the fixed number of connected
components of $U'\times_U U'$, there are at most $|P|^c$ possible
descent isomorphisms. The cocycle condition on the triple product
can only reduce this number. Finite-surjective proper descent of
finite etale schemes, together with full faithfulness for the
action morphisms, gives effective descent of these $P$-torsors
\citep[Expose~IX, Theorem~4.7]{C_SGA1}.

We may count all descended $P$-torsors, even if their monodromy
does not lie in $\operatorname{PSL}_r$; this only increases the
bound. Thus the gluing contributes at most $q^{cr^2}$ and proves
the displayed estimate. Distinct epimorphism kernels are separated
by projective conjugacy classes. The order estimate and the
generating-tuple argument in
Corollary~\ref{D:surface-bounded-centre-simple-product} now give
the stated finite generation of the simple-product quotient.
No degree-two cohomology assertion for the nonnormal group is
used in this proof.
\end{proof}
